Cette partie décrit l'architecture électronique mise en œuvre dans le cadre du projet. Elle présente l'ensemble des choix de conception réalisés afin d'assurer l'alimentation, le contrôle et l'actionnement de la maquette, tout en respectant les contraintes électriques et fonctionnelles identifiées.

L'architecture globale du système est d'abord introduite au moyen d'un schéma bloc, permettant de mettre en évidence les différents sous-ensembles et leurs interactions. Les solutions retenues pour la distribution d'énergie ainsi que pour l'entraînement des moteurs sont ensuite détaillées, en particulier au regard des contraintes de courant et de dimensionnement des composants.

Une attention particulière est portée au choix du microcontrôleur, élément central du système de commande, ainsi qu'à l'implémentation de l'électronique de puissance associée aux drivers moteurs. Enfin, la réalisation du schéma électrique complet, le prototypage sur breadboard et la conception du \gls{pcb} sont présentés et justifiés.
\section{Alimentation}
\section{Schéma bloc}
\section{Drivers moteurs}
\subsection{Courant max}
\section{Microcontrôleur}
\section{Schéma électrique}
\section{Breadboard}
\section{\gls{pcb}}


